\documentclass{beamer}

\usepackage[utf8]{inputenc}
\usepackage[spanish]{babel}
\usepackage[T1]{fontenc}
\usepackage{graphicx}
\usepackage{tikz}

\renewcommand\shorthandsspanish{}
\noextrasspanish

\title{Python}
\author{
Guillem Borrell i Nogueras\\
\texttt{guillemborrell@gmail.com}
}

\begin{document}

\begin{frame}
\begin{center}
 \includegraphics[width=9cm]{files/python-logo-generic.pdf}\\
 % python-logo-generic.pdf: 389x115 pixel, 72dpi, 13.72x4.06 cm, bb=0 0 389 115
\begin{large}
\textbf{Python en Supercomputación}
\end{large}\\
Charla Técnica\\

Guillem Borrell i Nogueras\\

ETSIA, Octubre 2007
\end{center}

\end{frame}

\begin{frame}
 \frametitle{Charla técnica}
\begin{itemize}
 \item Python Vs. Matlab.
 \item Más sobre los wrappers.
 \item Arrays en C, Fortran y Python.
 \item GIL.
 \item Uso de F2Py.
 \item ctypes.
 \item Python en paralelo.
\end{itemize}

\end{frame}

\begin{frame}
 \frametitle{Python vs. Matlab}
\begin{center}
 ¿Ayuda Matlab a manejar la complejidad?
\end{center}
\begin{itmeize}
 \item Cada función es un archivo
 \item No existe la modularidad
 \item Orientación a objetos inútil
 \item No existen los namespaces
 \item wrapper $eq$ mex
\end{itmeize}

\begin{frame}
\frametitle{Más sobre los wrappers}
\begin{itemize}
 \item Un wrapper es un adaptador entre los tipos de C, Fortran y los del lenguaje interpretado.
 \item En un lenguaje interpretado los tipos siempre son structs.
 \item En python los tipos tienen $\sim$ 30 elementos.
 \item El intérprete tiene su propio modelo de stack.
\end{itemize}
\end{frame}

\begin{frame}
 \frametitle{Arrays}
\begin{center}
 El objetivo es comunicar Arrays en C, Fortran con los arrays de \emph{numpy}.
\end{center}

\end{frame}



\end{document}
